\section{Tính ứng dụng của khoa học}

\

Quá trình phát triển của khoa học tất yếu sẽ dẫn đến sự phát triển của kĩ thuật và công nghệ. Cho tới hiện tại, chúng ta đã trải qua bốn cuộc cách mạng công nghiệp, mỗi cuộc cách mạng gắn liền với những phát minh vĩ đại của nhân loại.

Cách mạng công nghiệp lần thứ nhất gắn liền với phát minh ra động cơ hơi nước. Lần đầu tiên trong lịch sử, sức cơ bắp của con người và động vật được thay thế bằng sức mạnh của máy móc, cụ thể là những thiết bị được thiết kế dựa trên sự chuyển hóa nhiệt năng và cơ năng.
\begin{itemize}
    \item Năm 1764, Giêm Ha-gri-vơ\footnote{James Hargreaves (1720 -- 1778)} phát minh ra máy kéo sợi. Đến cuối thế kỉ \Rmnum{18}, hơn $\num{20000}$ máy kéo sợi như vậy được sử dụng tại Anh.
    \item Mặc dù Giêm Oát\footnote{James Watt (1736 -- 1819).} hay được coi như cha đẻ của nền công nghiệp lần thứ nhất, ông không phải là người đầu tiên phát minh ra động cơ hơi nước. Thực tế, ông dựa trên những nguyên mẫu sơ khai của Niu-ca-mân\footnote{Thomas Newcomen (1664 -- 1729).} và cải tiến nó, giúp động cơ hơi nước hoạt động hiệu quả hơn.
    \item Hen-ri Cót\footnote{Henry Cort (1740 -- 1800).} vào năm 1784 đã phát minh ra phương pháp mới trong luyện kim, cho phép sản xuất sắt thép chất lượng cao với số lượng lớn.
    \item Năm 1804, Ri-sát Tre-vi-thích\footnote{Richard Trevithick (1771 -- 1833).} và sau đó là Gioóc-giơ Sti-phen-xơn\footnote{George Stephenson (1781 -- 1848).} vào năm 1814 đã đặt đầu máy hơi nước chạy trên đường ray. Đây là hình ảnh biểu tượng nhất mà có thể bạn đọc thường được thấy ở các bảo tàng. Còn trong trường hợp bạn đọc chưa từng thấy bao giờ, hãy xem hình \ref{fig:khoa_hoc:gioi_thieu_ngan:tinh_ung_dung:tau_hoi_nuoc_co_dien}.
\end{itemize}

\begin{figure}[ht]
    \centering
    \includegraphics[width=0.6\textwidth]{images/steam_engine-masoodaslami.jpg}
    \caption{Bánh xe của một đoàn tàu hơi nước cổ điển, biểu tượng của cách mạng công nghiệp lần thứ nhất. \\ Nguồn: https://www.pexels.com/photo/close-up-of-historic-steam-locomotive-wheels-29831512/}
    \label{fig:khoa_hoc:gioi_thieu_ngan:tinh_ung_dung:tau_hoi_nuoc_co_dien}
\end{figure}

Cách mạng công nghiệp lần thứ hai bắt đầu từ nửa sau của thế kỉ \Rmnum{19} cho đến trước Thế chiến thứ nhất. Đây là kỉ nguyên của điện khí hóa, sản xuất hàng loạt và những bước tiến trong lĩnh vực truyền thông. Có thể kể đến một vài thành tựu như:

\begin{itemize}
    \item Năm 1865, Mắc-xoen\footnote{James Clerk Maxwell (1831--1879).} đã đưa ra dự đoán về trường điện từ. Dựa vào nó, Héc\footnote{Heinrich Hertz (1857 -- 1894).} đã thực hiện các thí nghiệm vào năm 1887 để tạo ra và phát hiện sóng vô tuyến, chứng minh rằng những gì Mắc-xoen dự đoán là hoàn toàn chính xác. Từ đó, bắt đầu ra đời ra-đi-ô và các thiết bị điện tử.
    \item Vào năm 1869, Men-đê-lê-ép\footnote{Dmitri Ivanovich Mendeleev (1834 -- 1907).} đã tạo ra một hệ thống phân loại các nguyên tố hóa học dựa trên khối lượng nguyên tử. Qua kết quả này, chúng ta có thể có cái nhìn tổng quát về tính chất của các nguyên tố hóa học và dự đoán được sự tồn tại của các nguyên tố chưa được phát hiện. Để vinh danh cho sự phát hiện vĩ đại này, nguyên tố men-đê-li-vi (Md) với số nguyên tử \num{101} được đặt tên theo ông.
          \begin{figure}[ht]
              \centering
              \includegraphics[width=0.6\textwidth]{images/bang_tuan_hoan_hoa_hoc-polina-tankilevitch.jpg}
              \caption{Một phần của bảng tuần hoàn hóa học hiện đại. \\ Nguồn: https://www.pexels.com/photo/photo-of-periodic-table-of-elements-3735778/}
              \label{fig:khoa_hoc:gioi_thieu_ngan:tinh_ung_dung:bang_tuan_hoan_hoa_hoc}
          \end{figure}
    \item Động cơ được chuyển từ chạy bằng hơi nước sang động cơ đốt trong. Ốt-tô\footnote{Nikolaus Otto (1832 -- 1891).} đã phát minh ra động cơ đốt trong bốn thì (hút - nén - nổ - xả) vào năm 1876. Nền tảng lí thuyết này lần đầu được đưa vào thực tế vào năm 1886 với việc ô tô ba bánh chạy bằng xăng ra đời bởi Các Ben-xơ\footnote{Karl Benz (1844 -- 1929).}. Không chạy bằng xăng thì chạy bằng đi-ê-den, với cơ chế cháy nổ do chính Đi-ê-den\footnote{Rudolf Diesel (1858 -- 1913). Giờ bạn đọc đã biết nguồn gốc tên gọi loại nhiên liệu này từ đâu ra rồi đó!} đề xuất. Thay vì sử dụng tia lửa điện như động cơ chạy bằng xăng, động cơ đi-ê-den sử dụng khí nén để làm nóng nhiên liệu đến mức tự bốc cháy, giúp tăng hiệu suất nhiệt đáng kể.
    \item Ngoài động cơ đốt trong, điện cũng ngày càng trở nên phổ biến hơn trong đời sống. Uy-li-am Xtơ-diên\footnote{William Sturgeon (1783 -- 1850).} vào năm 1832 đã sáng chế ra động cơ điện một chiều có sử dụng bộ cổ góp để đảo chiều dòng điện, giúp trục quay liên tục. Động cơ điện xoay chiều thì đến với cuộc chơi muộn hơn, khi Tét-xla\footnote{Nikola Tesla (1856 -- 1943).} nhận bằng sáng chế cho loại động cơ không đồng bộ vào năm 1888. Mặc dù đã tồn tại hơn một thế kỉ, các loại động cơ điện vẫn đóng một vai trò chủ chốt trong các ứng dụng công nghiệp và thương mại.
          \begin{figure}[ht]
              \centering
              \includegraphics[width=0.6\textwidth]{images/dong_co-wolsink.jpg}
              \caption{Động cơ điện. \\ Nguồn: https://www.pexels.com/photo/industrial-motor-on-wooden-surface-in-workshop-34194564/}
              \label{fig:khoa_hoc:gioi_thieu_ngan:tinh_ung_dung:dong_co_dien}
          \end{figure}
    \item Lí thuyết về năng lượng hạt nhân cũng được phát triển mạnh mẽ trong cách mạng công nghiệp lần thứ hai. Ma-ri Quy-ri\footnote{Marie Curie (1867 -- 1934).} cùng với chồng là Pi-e Quy-ri\footnote{Pierre Curie (1859 -- 1906).} đã đặt ra thuật ngữ ``phóng xạ'' và phát hiện các nguyên tố mới vào năm 1898, mở ra cánh cửa nhìn vào bên trong nguyên tử. Bảy năm sau đó, Anh-xtanh công bố phương trình nổi tiếng nhất thế giới $$E = mc^2,$$ chứng minh một lượng vật chất cực nhỏ có thể biến đổi thành một nguồn năng lượng khổng lồ. Từ đó, nền móng của công nghệ năng lượng hạt nhân được hình thành, với việc ứng dụng vào các nhà máy điện hạt nhân hiện đại ngày nay.
\end{itemize}

Bạn đọc đã có thể thấy một lượng thành tựu lớn hơn hẳn của nền công nghiệp lần thứ hai so với lần thứ nhất. Nếu tiếp tục liệt kê các thành tựu như vậy với các nền công nghiệp tiếp theo thì dung lượng của cả chương sẽ bằng cả một cuốn sách! Do đó, tác giả sẽ tóm tắt lại những kết quả nổi bật từ nền công nghiệp thứ ba tiến tới hiện tại.

Cách mạng công nghiệp lần thứ ba còn được gọi là cách mạng kĩ thuật số. Đây là khi nhân loại bắt đầu chuyển không gian làm việc từ tương tự sang số với sự ra đời của tran-xi-to vào cuối những năm 40 của thế kỉ \Rmnum{20}. Từ đó, các mạch tích hợp và những chiếc máy tính cá nhân dần dần đến được với công chúng. Hiện nay, gần như bất cứ mạch điện tử nào cũng đều chứa rất nhiều tran-xi-to bán dẫn. Hình \ref{fig:khoa_hoc:gioi_thieu_ngan:tinh_ung_dung:transistor} là một ví dụ, nhưng đừng lầm tưởng là mọi tran-xi-to cũng đều có kích cỡ giống như trong ví dụ. Một con vi xử lí hiện đại, mặc dù kích cỡ tổng thể chỉ khoảng vài xăng-ti-mét vuông, có thể có hàng tỉ tran-xi-to!

\begin{figure}[ht]
    \centering
    \includegraphics[width=0.6\textwidth]{images/transistor-tima-miroshnichenko.jpg}
    \caption{Tran-xi-to 15N03L, cùng với rất nhiều phần tử mạch điện tí hon khác. \\ Nguồn: https://www.pexels.com/photo/man-person-smartphone-industry-6755141/}
    \label{fig:khoa_hoc:gioi_thieu_ngan:tinh_ung_dung:transistor}
\end{figure}

Năm 1989, Tim Bơ-nơ Li\footnote{Tim Berners-Lee, sinh năm 1955.} đã phát minh ra mạng toàn cầu (World Wide Web --- www), cho phép trao đổi thông tin giữa con người nhanh chóng với khoảng cách lớn. Ngoài ra, đây là thời kì mà con người bắt đầu bước ra ngoài vũ trụ. Việt Nam cũng không nằm ngoài cuộc chơi này, với Phạm Tuân là người đầu tiên của Việt Nam được rời khỏi Trái Đất.

Chúng ta đang sống ở kỉ nguyên của công nghệ 4.0 (với nhiều nhà khoa học còn đang đưa ra luận điểm rằng chúng ta đã bước sang cách mạng lần thứ năm). Cách mạng này thường được gọi là kỉ nguyên của trí tuệ.
\begin{itemize}
    \item Về mảng vật lí, tự động hóa được đẩy lên đỉnh cao khiến cho rô-bốt ngày càng trở nên thông minh hơn và có thể hành động trên những phạm vi phức tạp hơn. Ngoài ra, điều này dẫn đến các sản phẩm có thể được in ra từ dữ liệu số, làm giảm thời gian thiết kế và chế tạo sản phẩm.
    \item Về mặt hóa học và vật liệu, chúng ta dần dần thành công trong việc giảm giá thành của các nguồn năng lượng tái tạo.
    \item Còn về sinh học, công nghệ gien, y học chính xác và nông nghiệp thông minh ngày càng trở nên quen thuộc với đời sống thường ngày.
\end{itemize}
Có thể nói chúng ta đã đạt đến một bước nhảy vọt mà có khi chưa từng một ai ở bất cứ thế kỉ trước tưởng tượng đến.\footnote{Sẽ còn tuyệt vời hơn nếu như mọi người trên thế giới biết phân phối nguồn lực và tài nguyên để đem đến cuộc sống đầy đủ và hạnh phúc cho mỗi dân tộc. Tuy nhiên, rào cản lớn nhất --- tư bản chủ nghĩa --- vẫn còn tồn tại.} Và đương nhiên, tác giả muốn nhấn mạnh rằng tất cả những thành tựu đều đến từ việc phát hiện và ứng dụng hợp lí các kết quả khoa học.